% (c) 2026 Stephanie Märklin, OST Ostschweizer Fachhochschule
%
% !TEX root = ../../buch.tex
% !TEX encoding = UTF-8
%
% Rolle im Paper: Realität in ein mathematisches Objekt übersetzen;
% Annahmen benennen
%
% Stichworte:
% 17.2 Das Flugzeug als dynamisches System (Vorspann, unnummeriert)
% - Ziel: Flugzeug -> dynamisches System übersetzen; Begriff einführen
%   (mathematisches Objekt, dessen Zustand sich mit der Zeit ändert)
% - Begriffe aus 17.1 beschreiben nur die Lage; Erweiterung um
%   Grössen, die die Bewegung erfassen
% - Beschränkung: Seitenbewegung (lateral motion) = Gieren, Schieben,
%   Rollen -- nur diese sind an der Dutch Roll beteiligt
% - Annahme: Störung klein, Betrachtung nahe der Ruhelage
% - Fahrplan am Ende ("Für diese Übersetzung ..."): Zustandsvektor ->
%   Bewegungsgesetz -> Systemmatrix

% 17.2.1 Der Zustandsvektor der Seitenbewegung
% - vier Grössen einzeln: Schiebewinkel beta (formalisiert das
%   Schieben aus 17.1.2), Rollrate p, Gierrate r, Rollwinkel phi
% - Tabelle: Symbol, Name, Bedeutung
% - Notation nach Nelson \cite (ein Satz, deckt die Tabelle ab)
% - Zustandsvektor x = [beta, p, r, phi]^T (equation, Label)
% - ein Satz Rückgriff: beta verbindet Gieren und Rollen;
%   Konsequenz folgt in Abschnitt 17.4 (Vorverweis, kein Vorgriff)
% - Schluss (Paar, keine Überleitung): x = Momentaufnahme zur Zeit t;
%   Ableitung xdot beschreibt, wie sich der Zustand verändert

% 17.2.2 Das Bewegungsgesetz
% - Andocken: wir suchen den Zusammenhang zwischen x und xdot
% - für kleine Störungen um die Ruhelage linear -> xdot = Ax
%   (equation, Label); \emph{Systemmatrix} benennen;
%   Herleitung: Nelson \cite
% - Gleichung lesen: Änderung jeder Zustandsgrösse ist eine
%   Linearkombination der momentanen Zustandsgrössen; A enthält
%   die Koeffizienten

% 17.2.3 Die Systemmatrix
% - Andocken: "Für die Seitenbewegung lautet die Systemmatrix ..."
% - A vollständig (equation, Label; exakt die Vortrags-Fassung)
% - neue Symbole erklären: g Erdbeschleunigung, V Fluggeschwindigkeit
% - offener Punkt: Schreibweise ggü. Nelson Gl. 5.33 vereinfacht --
%   Halbsatz ergänzen oder Konvention mit Formelskelett/MATLAB klären
% - zwei Zeilen vorlesen: pdot = Lbeta*beta + Lp*p + Lr*r (gehaltvoll)
%   und phidot = p (Buchhaltung; löst Rollen/Rollrate/Rollwinkel ein)
% - Einträge = aerodynamische Derivative; EIN Beispiel: Lbeta =
%   Rollmoment durch Schieben + Böen-Analogie (noch von Stephanie)
% - Struktur fix, Werte flugzeug-/flugzustandsspezifisch
%   (Tabelliert-Halbsatz gestrichen)
% - Schluss (Feststellung): "Damit ist das Flugzeug in ein
%   dynamisches System übersetzt."


\section{Das Flugzeug als dynamisches System
\label{aircraft:section:system}}
\kopfrechts{Dynamisches System}
Um die Bewegung nach einer Störung berechnen zu können, übersetzen wir
das Flugzeug in diesem Abschnitt in ein \emph{dynamisches System}.
Ein mathematisches Objekt, dessen Zustand sich mit der Zeit verändert.
Dazu führen wir zuerst den Zustandsvektor ein, dann das
Bewegungsgesetz, das seine Änderung bestimmt, und die Systemmatrix mit ihren
Einträgen.

Wir beschränken uns dabei auf die \emph{Seitenbewegung} (lateral
motion): das Gieren, das Schieben und das Rollen.
Wir nehmen an, dass die Störung klein ist, und betrachten die
Bewegung in der Nähe der Ruhelage.

\subsection{Der Zustandsvektor der Seitenbewegung
\label{aircraft:subsection:zustandsvektor}}
Der Zustand ist die Momentaufnahme der Bewegung.
Für die Seitenbewegung genügen dazu vier Grössen.
Der \emph{Schiebewinkel} $\beta$ formalisiert das Schieben aus
Abschnitt~\ref{aircraft:subsection:dutchroll}:
Er misst den Winkel zwischen der Längsachse und der tatsächlichen
Bewegungsrichtung des Flugzeugs.
Die \emph{Rollrate} $p$ und die \emph{Gierrate} $r$ sind die
Änderungsraten des Rollens und des Gierens.
Der \emph{Rollwinkel} $\varphi$ misst die Drehung um die Längsachse
gegenüber der Ruhelage.
Wir folgen dabei der Notation von Nelson~\cite{aircraft:nelson}, der
die Seitenbewegung mit denselben Grössen beschreibt.

Wir fassen die vier Grössen zum \emph{Zustandsvektor}
\begin{equation}
    \vec{x}(t)
    =
    \begin{pmatrix}
        \beta(t) \\ p(t) \\ r(t) \\ \varphi(t)
    \end{pmatrix}
    \label{aircraft:eqn:zustandsvektor}
\end{equation}
zusammen.
Der Zustandsvektor ist die Momentaufnahme der Bewegung zur Zeit $t$.
Seine Ableitung $\dot{\vec{x}}(t)$ beschreibt, wie sich der Zustand
verändert.

Der Schiebewinkel verbindet dabei das Gieren mit dem Rollen.
Welche Folgen diese Verbindung für die Berechnung hat, zeigt sich in
Abschnitt~\ref{aircraft:subsection:eigenwertgleichung}.
% // TODO Referenz anpassen oder evtl. ganz streichen

\subsection{Das Bewegungsgesetz
\label{aircraft:subsection:bewegungsgesetz}}
Wir suchen den Zusammenhang zwischen dem Zustand $\vec{x}$ und seiner
Änderung $\dot{\vec{x}}$.
Für kleine Störungen um die Ruhelage ist dieser Zusammenhang linear
und es gilt
\begin{equation}
    \dot{\vec{x}}
    =
    A\vec{x}
    \label{aircraft:eqn:bewegungsgesetz}
\end{equation}
mit einer Matrix $A$, der \emph{Systemmatrix}.
Die Herleitung findet sich bei Nelson~\cite{aircraft:nelson}.
Die Änderung jeder Zustandsgrösse ist also eine Linearkombination
der momentanen Zustandsgrössen.

\subsection{Die Systemmatrix und Ihre Einträge
\label{aircraft:subsection:systemmatrix}}
Für die Seitenbewegung lautet die Systemmatrix
\begin{equation}
    A
    =
    \begin{pmatrix}
        Y_\beta & Y_p & Y_r & g/V \\
        L_\beta & L_p & L_r & 0   \\
        N_\beta & N_p & N_r & 0   \\
        0       & 1   & 0   & 0
    \end{pmatrix}.
    \label{aircraft:eqn:systemmatrix}
\end{equation}
Ausgeschrieben ist die zweite Zeile der Systemmatrix
~\eqref{aircraft:eqn:bewegungsgesetz} $\dot{p} = L_\beta\beta + L_p p + L_r r$.
Die Änderung der Rollrate entsteht aus dem momentanen Schiebewinkel,
der momentanen Rollrate und der momentanen Gierrate.
Die Einträge der Systemmatrix heissen aerodynamische Derivative.
Das Derivativ $L_\beta$ gibt zum Beispiel an, wie stark der
Schiebewinkel auf das Rollen wirkt.
Die Struktur der Matrix ist für alle Flugzeuge gleich.
Die Werte der Derivative hängen vom Flugzeug und vom Flugzustand ab.